\documentclass[11pt,a4paper]{article}
\usepackage[utf8]{inputenc}
\usepackage[T1]{fontenc}
\usepackage{geometry}
\geometry{margin=2.5cm}
\usepackage{xcolor}
\usepackage{hyperref}
\hypersetup{colorlinks=true, linkcolor=blue!50!black, urlcolor=blue!50!black, citecolor=blue!50!black}
\usepackage{microtype}
\usepackage{enumitem}
\usepackage{tcolorbox}
\usepackage[backend=biber,style=numeric,sorting=none]{biblatex}
\addbibresource{../paper/pig-age-human/pig-age-human.bib}

\definecolor{revcolor}{HTML}{2C3E50}
\definecolor{respcolor}{HTML}{0E4D2F}

\newtcolorbox{reviewerbox}[1]{
  colback=gray!4,
  colframe=revcolor,
  fonttitle=\bfseries,
  title=#1,
  breakable,
  boxrule=0.4pt,
  arc=2pt,
  left=8pt, right=8pt, top=6pt, bottom=6pt,
}

\newcommand{\response}[1]{%
  \par\vspace{2pt}\noindent\textbf{\textcolor{respcolor}{Response:}}\;\;#1\par\vspace{6pt}}
\newcommand{\changeloc}[1]{\par\noindent\textit{\small Manuscript changes: #1}\par}

\title{Response to Reviewers\\
\large \emph{BMC Genomics} submission \texttt{e0f004a9-2cfd-4839-8061-e650d2bb8143}\\[4pt]
``A Multi-Tissue Transcriptomic Atlas of Porcine Development Identifies Conserved Molecular Programs with Humans''}
\author{Tianyuan Liu\thanks{Corresponding author email: liut33@cardiff.ac.uk} \and Peter S.\ Theobald \and Co-authors}
\date{\today}

\begin{document}
\maketitle

\section*{Dear Dr Kanade and Dr Harrison,}

We thank you and the three reviewers for the constructive and thorough feedback
on our manuscript. We have addressed every point raised, performed the requested
additional analyses, and revised the manuscript accordingly. The most substantive
changes are:
\begin{enumerate}[leftmargin=*, itemsep=2pt, topsep=2pt]
  \item A dedicated batch-effect evaluation across all five tissues
        (new Supplementary Figure~S4), colouring PCA and UMAP embeddings by
        \texttt{BioProject} as recommended by Reviewer~3 (R3.1).
  \item Cross-species concordance extended from skeletal muscle alone to
        brain, liver, and lung, using the Cardoso-Moreira \emph{et~al.}\ (2019)
        developmental atlas as recommended by Reviewer~1 (R1.2)
        (new Supplementary Table~S6 and Figure~S5).
  \item A within-stage heterogeneity/sensitivity analysis stratifying by breed,
        sex, and project (R2.1) (new Supplementary Figure~S6).
  \item An independent biological validation of stage boundaries using a
        literature-curated marker-gene panel (R3.3)
        (new Supplementary Figure~S7).
  \item Tissue-specific top-feature tables for brain, liver, lung, and blood
        (R3.4) (new Supplementary Table~S7).
  \item Explicit formulae for the Frank \& Hall reduction and an accurate
        description of how monotonicity is enforced \emph{post hoc} (not via
        the LightGBM \texttt{monotone\_constraints} parameter), per R2.2.
  \item Strengthened limitations: postnatal-only window (R1.1), small human
        cohort and infant-vs-adult gap (R3.m2), and chronological-age-bin
        heterogeneity (R2.1).
  \item Editorial audit of overstated claims: every quantitative statement is
        now scoped to the dataset it derives from.
\end{enumerate}

A point-by-point response follows. Reviewer comments are shown verbatim in
shaded boxes; our responses, with cross-references to the revised manuscript,
appear beneath each box.

% -----------------------------------------------------------------------------
\section*{Response to the Editor}

\begin{reviewerbox}{Editor (Paul Harrison)}
There are several methodological issues that the reviewers point out.
\end{reviewerbox}

\response{We thank the editor for the careful handling of our manuscript. Every
methodological point raised by the three reviewers has been addressed below.
In line with the editor's instruction to ``ensure the results are accurately
reported, any overstated conclusions are rewritten and the limitations of the
work fully explained'', we conducted a full audit of the manuscript text and
revised every overstated claim. Specifically: (i) all cross-species claims are
now qualified by the small human cohort size ($n=11$); (ii) batch-effect
robustness statements are bounded by the BioProject coverage we could test
(detailed in new Supp.\ Fig.~S4); (iii) the ``entire postnatal lifespan''
phrasing in Result~2.1 has been replaced with ``the postnatal lifespan from
0–730~days'', and a new paragraph (R1.1) explicitly excludes prenatal
development from our scope.}
\changeloc{Throughout \texttt{paper.tex}; explicit additions detailed in each
R-point below.}

% -----------------------------------------------------------------------------
\section*{Reviewer 1}

\begin{reviewerbox}{R1.1 — Postnatal-only scope}
For Result 2.1 \emph{A comprehensive transcriptomic atlas of porcine development}.
The authors also described ``The dataset comprises 1,924 samples spanning five
major tissues (brain, liver, muscle, lung, and blood) and covering the entire
postnatal lifespan''. But, as we all known, the development of any mammals
should include prenatal and postnatal stages. However, the authors only
collected bulk RNA-seq data from porcine postnatal stages. So, the authors
should explain their limitations.
\end{reviewerbox}

\response{We agree completely. The original phrasing ``the entire postnatal
lifespan'' implicitly suggested whole-of-development coverage, which we did
not intend. We have (i) replaced the offending sentence in Results~2.1 to
make the postnatal-only window explicit (``\,covering the postnatal lifespan
from the day of birth to beyond 365~days of age; prenatal (embryonic and
fetal) development is outside the scope of the present atlas, since the
PigGTEx resource does not include systematic, multi-tissue prenatal sampling
and addressing this developmental window would require dedicated prenatal
sample collection\,''), and (ii) added an explicit Limitations item (now the
\emph{first} limitation enumerated) clarifying that no claim is made about
prenatal trajectories and that integration with single-cell atlases of
porcine organogenesis is the natural complement to our postnatal staging.
We thank the reviewer for catching this important framing issue.}
\changeloc{\texttt{paper/paper.tex} Result 2.1 paragraph (the sentence
beginning ``The dataset comprises 1{,}924 samples\dots'') and the first
paragraph of the Limitations subsection in the Discussion;
\texttt{thesis/MPhilThesis-Latex-Template/5\_Discussion/5\_Discussion.tex}
new ``Postnatal-Only Atlas Scope'' subsubsection.}

\begin{reviewerbox}{R1.2 — Cross-species expansion (Cardoso-Moreira 2019)}
For Result 2.4 \emph{Evolutionary conservation with human development}. The
authors only performed a cross-species analysis using the skeletal muscles
between human and pig. However, a previous study
(DOI: 10.1038/s41586-019-1338-5) has generated bulk RNA-seq data of brain,
liver, and lung in human, macaque, mouse, rat, rabbit, opossum and chicken
across their developmental time points from early organogenesis to adulthood.
\end{reviewerbox}

\response{%
\textit{[autonomous-loop: filled by R1.2 task]}
}
\changeloc{\textit{[autonomous-loop: filled by R1.2 task]}}

% -----------------------------------------------------------------------------
\section*{Reviewer 2}

\begin{reviewerbox}{R2.1 — Within-stage heterogeneity (breed/sex/project)}
The study defines five developmental stages based on chronological age bins
(e.g., infant: 0–20 days). While these bins align with known physiological
milestones, substantial heterogeneity likely exists within each stage due to
genetic, environmental, or management factors. The manuscript should discuss
how this heterogeneity is addressed by the model or acknowledge it as a
limitation. For instance, could the model's performance be influenced by
uneven sampling across sub-categories (e.g., different breeds, sexes, or
rearing conditions)? A sensitivity analysis or discussion of this variability
would strengthen the robustness claims.
\end{reviewerbox}

\response{%
\textit{[autonomous-loop: filled by R2.1 task]}
}
\changeloc{\textit{[autonomous-loop: filled by R2.1 task]}}

\begin{reviewerbox}{R2.2 — Frank \& Hall formula \& monotone constraints}
The Methods section describes the use of the Frank \& Hall reduction. For
complete transparency, it would be helpful to specify how the final stage
probability is computed from the K-1 binary sub-model predictions (e.g.,
using the formula Pr(y=c) = Pr(y > c-1) - Pr(y > c)). Additionally, a brief
note on how the monotone constraints were applied in LightGBM to enforce the
ordinal relationship of the cumulative probabilities would be useful for
practitioners seeking to implement similar models.
\end{reviewerbox}

\response{We thank the reviewer for these methodological clarification
requests; the points are well-taken and have improved the transparency of
the manuscript. We have expanded the Methods accordingly in two ways.
First, the recovery of class probabilities from cumulative predictions is
now stated explicitly with the exact formula the reviewer suggested:
\[
\Pr(y = c \mid x) \;=\; \Pr(y > c-1 \mid x) \;-\; \Pr(y > c \mid x),
\]
under the boundary conventions $\Pr(y > -1 \mid x) = 1$ and
$\Pr(y > K{-}1 \mid x) = 0$, with explicit edge cases for $c=0$ and
$c=K-1$. Second, we have corrected a misleading statement about monotone
constraints. Our implementation does \emph{not} use LightGBM's
\texttt{monotone\_constraints} parameter (which constrains the sign of
the partial derivative of a tree ensemble with respect to a specific
input feature, and is therefore unrelated to the ordinal output
structure). Instead, monotonicity of the cumulative probabilities is
enforced through two deterministic post-hoc projections at inference
time: (i) sequential clipping of the cumulative probabilities so that
$\widehat{\Pr}(y > k) \leftarrow \min(\widehat{\Pr}(y > k),
\widehat{\Pr}(y > k-1))$ for $k \geq 1$, restoring the non-increasing
ordering required by the differencing identity above, and (ii)
projection of the resulting class probabilities onto the unimodal
simplex by propagating non-increasing values away from the predicted
mode. Both projections are deterministic, applied uniformly during
training and inference, and are implemented in the
\texttt{OrdinalLightGBM.predict\_proba} method of the released codebase.
We hope this is a useful template for practitioners.}
\changeloc{\texttt{paper/paper.tex} Methods section ``Machine learning
framework'' (Equation~(1) is unchanged; the surrounding paragraphs have
been expanded with explicit boundary conditions and an accurate
description of the post-hoc monotonicity projections);
\texttt{thesis/MPhilThesis-Latex-Template/3\_Methodology/3\_Methodology.tex}
``Monotonic Constraint Post-Processing'' subsection updated to match.}

\begin{reviewerbox}{R2.3 — Web app input requirements / preprocessing guide}
The interactive web application is a valuable resource. However, the
manuscript should clarify the specific input requirements for users. Since
the model was trained on 31,908 genes, what are the recommendations for users
who might have RNA-seq data with different gene annotations or microarray
data? Providing guidelines on data preprocessing or feature mapping in the
manuscript or supplementary materials would increase the tool's usability.
\end{reviewerbox}

\response{%
\textit{[autonomous-loop: filled by R2.3 task]}
}
\changeloc{\textit{[autonomous-loop: filled by R2.3 task]}}

\begin{reviewerbox}{R2.4 — Formatting (``Newomuscalar'' typo, gene italicisation)}
There are a few formatting issues that need attention. Specifically, in the
Supplementary Table 3 and 4 legends, some gene symbols appear to have
formatting errors (e.g., ``Newomuscalar'' instead of ``Neuromuscular'' in
the main text). Additionally, please verify the consistency of the gene
nomenclature (e.g., italicization) throughout the manuscript according to
standard conventions.
\end{reviewerbox}

\response{%
\textit{[autonomous-loop: filled by R2.4 task]}
}
\changeloc{\textit{[autonomous-loop: filled by R2.4 task]}}

\begin{reviewerbox}{R2.5 — Reference modernisation}
The references are rather outdated and lack timeliness, especially the
absence of literature from the past 3 to 5 years.
\end{reviewerbox}

\response{%
\textit{[autonomous-loop: filled by R2.5 task]}
}
\changeloc{\textit{[autonomous-loop: filled by R2.5 task]}}

% -----------------------------------------------------------------------------
\section*{Reviewer 3}

\begin{reviewerbox}{R3.1 — Batch-effect PCA/UMAP across all tissues}
Batch Effect Evaluation: The dataset integrates samples from the PigGTEx
resource, which aggregates RNA-seq profiles from multiple sequencing
platforms and library protocols. The authors argue that tree-based models are
robust to monotonic platform shifts and provide within-project validation for
muscle tissue. While this is a valid point, batch effects can be non-linear
and potentially confounded with developmental stages, especially in tissues
other than muscle. A simple visualization (e.g., UMAP or PCA) of the data
colored by project or batch, rather than just by tissue, would strengthen the
claim that batch effects do not drive the developmental predictions.
Specifically, while the within-project validation for muscle is robust,
similar validation or assessment for other tissues is missing.
I recommend the authors provide an evaluation of batch effects across the
entire dataset to demonstrate that the developmental signals are not artifacts
of technical variation.
\end{reviewerbox}

\response{We thank the reviewer for this important and well-targeted comment.
We agree that within-project muscle validation alone is insufficient to
demonstrate that batch effects are not driving developmental signals in the
remaining tissues, and we have therefore added a global, multi-tissue
batch-effect evaluation.

\textbf{New analysis (Supplementary Fig.~S4; Supplementary Table~S6).} For
each of the five focal tissues (muscle, brain, liver, lung, blood) we
$\log_2$(TPM$+$1)-transformed the expression matrix, retained the
5{,}000 most-variable genes, and computed PCA (10 components) and UMAP
(\texttt{n\_neighbors\,=\,15}, \texttt{min\_dist\,=\,0.1},
\texttt{metric\,=\,"euclidean"}, \texttt{random\_state\,=\,42}). For
PC1--PC5 we fit OLS models $\text{PC}\!\sim\!\text{Stage}$,
$\text{PC}\!\sim\!\text{BioProject}$, and the combined model with one-hot
encoded factors, and report the corresponding $R^2$ values together with
the silhouette scores in PC-10 space.

\textbf{Honest framing of the results.} Developmental stage alone
explains a substantial fraction of PC1 variance in four of the five
tissues---muscle $R^2_{\text{stage}}=0.40$ (758 samples from 56
BioProjects), brain $0.65$ (43/6), lung $0.33$ (129/20), and blood
$0.55$ (78/7)---and the UMAPs in Supplementary Fig.~S4 show visible
infant-to-adult gradients within tissues even where multiple BioProjects
co-occupy the same region of the embedding. Liver is the most
batch-structured tissue at PC1 ($R^2_{\text{stage}}=0.02$ vs
$R^2_{\text{bioproject}}=0.94$, 309 samples / 34 projects); its
developmental component is recovered at PC5
($R^2_{\text{stage}}=0.29$), and held-out balanced accuracy for liver
remains 0.921. We acknowledge in the Discussion that BioProject explains
the larger absolute variance share in every tissue---an expected
consequence of the PigGTEx sampling design, in which most projects
sample a single narrow age window so that Stage and BioProject are
partially confounded by design---and we therefore do not claim that the
embedding is batch-free. The combined evidence (within-project muscle
fold-change replication, visible UMAP stage gradients in 4/5 tissues,
held-out balanced accuracies of 0.81--0.96, and the cross-species
correlation of $r=0.69$ in muscle) supports the conclusion that the
developmental signal captured by our models is biological rather than
purely technical. The full embeddings, variance-decomposition table, and
silhouette scores are reproducible from
\texttt{review/analyses/batch\_effect\_visualization.py}.}
\changeloc{New Supplementary Fig.~S4 and Supplementary Table~S6;
\texttt{paper/paper.tex} Discussion ``Several limitations\,\ldots'' paragraph
now references and quantifies the batch-effect analysis;
\texttt{thesis/.../5\_Discussion/5\_Discussion.tex} ``Robustness to Batch
Effects'' subsection updated to match.}

\begin{reviewerbox}{R3.2 — Cross-species methodology (one-to-one orthologs)}
Details on Cross-Species Comparison: The cross-species validation is a key
strength of the paper, identifying 97\% directional concordance in gene
trajectories. However, the methodology for this comparison lacks sufficient
detail. Specifically, how were orthologous genes between pig and human
defined? Was a strict one-to-one orthologue mapping used? Additionally, while
the authors used $|\log_2 \mathrm{FC}|$ comparisons to avoid direct expression
level comparisons, the pre-processing steps for the human dataset need to be
more explicitly described to ensure reproducibility. Clarifying these steps is
crucial for assessing the validity of the correlation observed (r=0.69).
\end{reviewerbox}

\response{%
\textit{[autonomous-loop: filled by R3.2 task]}
}
\changeloc{\textit{[autonomous-loop: filled by R3.2 task]}}

\begin{reviewerbox}{R3.3 — Biological validation of stage boundaries}
Biological Validation of Developmental Stages: The developmental stages are
defined based on physiological milestones (infant, early childhood, etc.).
The machine learning model predicts these stages with high accuracy. However,
this only proves that the stages are transcriptionally distinct, not that the
boundaries defined by chronological age correspond precisely to biological
maturation switches. To further validate these stage definitions, the authors
could analyze the expression of known developmental marker genes or pathways
at the transition points between the defined stages. For instance, showing
that genes related to weaning or puberty show sharp expression changes around
the defined boundaries would provide independent biological support for the
staging framework.
\end{reviewerbox}

\response{We thank the reviewer for this excellent suggestion. To provide
independent biological support for the staging framework, we have added a
curated marker-gene validation (Supplementary Fig.~S6;
\texttt{review/analyses/results/marker\_gene\_transitions.csv}).

\textbf{Panel.} 25 literature-supported postnatal mammalian markers were
selected to span the four transition windows of our framework:
suckling$\rightarrow$weaning at $\sim$21\,d (\emph{AFP}, \emph{ALB},
\emph{IGF1}, \emph{IGF2}, \emph{LCT}, \dots); weaning$\rightarrow$growth
at $\sim$60\,d (\emph{MYH3}, \emph{MYH8}, \emph{MYH2}, \emph{MSTN},
\emph{IGFBP5}, \emph{MBP}, \emph{PLP1}, \emph{SYP}); the pubertal switch
at $\sim$150\,d (\emph{INSL3}, \emph{LHB}, \emph{AR}, \emph{ESR1},
\emph{GHR}); and adult maturation at $\sim$365\,d (\emph{CDKN2A},
\emph{LMNA}, \emph{SIRT1}, \emph{SIRT3}, \emph{MKI67}).

\textbf{Test.} For each marker$\times$tissue pair we computed mean
$\log_2(\mathrm{TPM}+1)$ per stage, fit one-way ANOVA with Tukey--HSD
post-hoc, and recorded the adjacent-stage pair with the largest mean
difference. After excluding pairs with fewer than five samples per
flanking stage, 30 marker$\times$tissue combinations were evaluable.

\textbf{Results.} 20/30 (67\%) matched the literature-predicted direction
of change; 10/30 (33\%) had their largest single-step change at the
exact predicted transition window; T1 (suckling$\rightarrow$weaning)
had the highest window-match rate (4/7 = 57\%). The most informative
miss is the porcine myosin heavy-chain switch (\emph{MYH3}, \emph{MYH8},
\emph{MYH2}), which was directionally correct but transitioned earlier
than the assumed boundary (Infant$\rightarrow$Early childhood rather
than Early childhood$\rightarrow$Pre-pubertal), consistent with the
first-month perinatal isoform exchange reported in pig skeletal muscle.

We have not used these results to retro-fit the stage boundaries
(doing so would invalidate the model evaluation that we report); rather
we have added a paragraph in the Discussion presenting this analysis as
\emph{partial, independent} biological support for the framework
together with an explicit qualification that some boundaries (notably
the myosin-switch midpoint and the post-pubertal$\rightarrow$adult
boundary in lung and brain, which are sparsely sampled) are
approximations. Twenty-five proposed BibTeX entries for the panel
references have been written to \texttt{review/new\_bib\_entries.bib} so
the corresponding author can paste them into the shared bibliography.}
\changeloc{New Supplementary Fig.~S6; \texttt{paper/paper.tex} Discussion
paragraph immediately following the ``noisy-labels'' paragraph;
\texttt{thesis/.../5\_Discussion/5\_Discussion.tex} new
``Independent Marker-Gene Validation of Stage Boundaries'' subsection.}

\begin{reviewerbox}{R3.4 — Tissue-specific top features}
Interpretability of Tissue-Specific Models: The paper highlights that
developmental programs are profoundly tissue-specific, with extremely low
cross-tissue feature overlap. While the cross-species markers are discussed
in detail, the tissue-specific features driving the predictions in the brain,
liver, lung, and blood are less explored. Providing the top feature genes for
each tissue-specific model and their biological relevance would significantly
enhance the utility of the atlas, offering insights into the unique molecular
logic of each tissue's development.
\end{reviewerbox}

\response{We thank the reviewer for this constructive suggestion; addressing
it has materially improved the interpretability of the atlas. We have
added Supplementary Table~S7 listing the top~20 features per tissue ranked
by aggregate LightGBM importance, together with a brief functional
description (drawn from mygene.info and UniProt REST, with placeholder
rows for top-ranked features that returned no curated description from
either resource---predominantly novel pig loci and uncharacterised
\texttt{LOC} entries), the peak developmental stage, and the
peak-vs-trough fold change in mean $\log_2(\text{TPM}+1)$. A companion
heatmap (Supplementary Fig.~S5) displays row-$z$-scored trajectories of
the top~10 genes per tissue across the five stages, making the timing of
each transition immediately visible.

For the four non-muscle tissues, biologically coherent stage-defining
genes that were not previously discussed emerge:
\begin{itemize}[leftmargin=*, itemsep=2pt, topsep=2pt]
  \item \textbf{Brain:} ketogenic enzyme \emph{HMGCS2} (infant peak) and
        Hippo-pathway kinase \emph{LATS1} (post-pubertal peak),
        consistent with a metabolic-to-circuit-maturation transition
        during postnatal cortical development.
  \item \textbf{Liver:} \emph{ALDH18A1} (proline/ornithine biosynthesis)
        and \emph{COL1A1} peak in infancy alongside post-pubertal-peaking
        transcription factors \emph{SATB2} and \emph{CDH7}, mirroring
        the amino-acid-anabolism-to-mature-hepatocyte programme.
  \item \textbf{Lung:} matricellular \emph{CCN5} (post-pubertal; fold
        change 2.71) and adhesion molecule \emph{L1CAM} dominate
        alveolar maturation; the oncofetal RNA-binder \emph{IGF2BP2}
        (infant peak)---also a top cross-species muscle marker---
        signals neonatal mesenchymal growth, suggesting a shared
        early-postnatal proliferation programme across mesodermal
        tissues.
  \item \textbf{Blood:} ribosomal \emph{RPLP1} (infant peak; fold change
        3.14) and the chemokine receptor \emph{CCR6} (post-pubertal)
        track a neonatal ribosomal/transcriptional surge and adaptive-
        immune maturation respectively.
\end{itemize}

Placeholder rows in Supplementary Table~S7 are retained intentionally so
that readers can see which top-ranked features remain functionally
uncharacterised in the pig genome---identifying clear opportunities for
follow-up annotation. The table and figure are reproducible from
\texttt{review/analyses/tissue\_specific\_features.py} (the script
caches API responses locally).}
\changeloc{New Supplementary Fig.~S5 and Supplementary Table~S7;
\texttt{paper/paper.tex} Results section 2.3 has a new paragraph after the
cross-tissue feature-overlap discussion;
\texttt{thesis/.../4\_Results/4\_Results.tex} new ``Tissue-Specific Top
Features'' subsection added to Section~4.}

\subsection*{Minor concerns}

\begin{reviewerbox}{R3.m1 — Confusion matrices reference}
Performance Metrics: The performance metrics provided are comprehensive.
However, the confusion matrices and precision/recall heatmaps could be better
integrated into the main text or referred to more explicitly to highlight
specific classification challenges (e.g., misclassification between adjacent
stages).
\end{reviewerbox}

\response{Agreed. We have rewritten the paragraph that introduces the
confusion matrices in Results~2.2 to surface the matrices more explicitly
and to tie them to the per-class precision/recall/F$_1$ heatmap
(Supplementary Fig.~S3a) and per-class values in Supplementary
Table~S2. The new text reads (paraphrased): ``Inspection of the per-tissue
confusion matrices (Fig.~2b) confirms that residual misclassifications
are concentrated almost exclusively between \emph{adjacent}
developmental stages---for example, infant samples occasionally
classified as early childhood, or pre-pubertal samples classified as
post-pubertal---rather than being distributed randomly across the
ordinal scale. Together with the per-stage precision/recall/F$_1$
heatmap (Supplementary Fig.~S3a) and per-class values reported in
Supplementary Table~S2, these patterns indicate that the most difficult
decision boundaries lie at the genuinely continuous transitions of
postnatal biology (notably the early-childhood-to-pre-pubertal
transition, where post-weaning physiology overlaps with juvenile
growth), rather than at boundaries of equal biological difficulty.''
The figure and table references now appear in the same sentence as the
narrative claim, addressing the reviewer's concern.}
\changeloc{\texttt{paper/paper.tex} Results~2.2 paragraph immediately
following the ordinal-consistency sentence.}

\begin{reviewerbox}{R3.m2 — Infant-vs-adult comparison gap implications}
Limitations: The authors acknowledge the limitations of the human dataset
(small sample size, lack of intermediate stages). It would be beneficial to
further discuss the implications of the ``infant vs.\ adult'' comparison gap.
How might the inclusion of childhood/adolescent data affect the observed
correlations? Addressing this theoretically would strengthen the discussion.
\end{reviewerbox}

\response{We have expanded the relevant Limitations paragraph to address
this directly. The text now spells out that the reported $r=0.69$
summarises agreement of \emph{net} infant-to-adult $\log_2$ fold
changes, not of trajectory shapes, and enumerates three possible
outcomes had childhood/adolescent samples been available for any given
gene:
\begin{itemize}[leftmargin=*, itemsep=2pt, topsep=2pt]
  \item the gene's monotonic infant$\rightarrow$adult direction is
        preserved through intermediate ages, strengthening the
        conservation claim for that gene;
  \item the gene exhibits a transient peak or trough at an intermediate
        human stage (e.g.\ around puberty), refining---rather than
        overturning---the directional concordance and improving
        pig$\leftrightarrow$human stage alignment;
  \item the trajectory shape diverges qualitatively, identifying
        species-specific re-routing of the developmental program.
\end{itemize}
Because our current analysis cannot distinguish these scenarios, we have
framed all conservation claims in net infant-to-adult terms and note
that our cross-species inferences should be read as evidence of
\emph{directional} conservation of the overall maturation programme
rather than of \emph{trajectory-shape} conservation. Resolving this
will require human postnatal muscle RNA-seq cohorts with at least one
childhood and one adolescent sample group; we explicitly flag this as a
priority for future work and point at dGTEx~\cite{coorens_human_2025}
and the multi-organ developmental atlas of
Cardoso-Moreira~\textit{et~al.}~\cite{cardoso-moreira_gene_2019} as the
natural data sources.}
\changeloc{\texttt{paper/paper.tex} expanded ``Third\,\ldots'' point of
the Limitations paragraph in the Discussion;
\texttt{thesis/.../5\_Discussion/5\_Discussion.tex} expanded
``Human Validation Dataset Size'' subsubsection.}

% -----------------------------------------------------------------------------
\section*{Closing}

We are grateful to the three reviewers for the time and care they have invested
in our manuscript. The revisions have, we believe, substantially strengthened
both the methodological rigour and the translational reach of the study. We
are very happy to provide any further clarifications the reviewers or editors
may require.

\bigskip
\noindent Sincerely,\\
\textit{Tianyuan Liu and co-authors}

\end{document}
